% Hexagon note H06b, 23 Nov 2026. CC0 1.0.
% Mathematics and Sections 3-8: GPT-6 Astra (ChatGPT Pro mode), November 2026.
% Sections 1-2, selection and closing: Claude Fable 5.2 at effort max, as editor.
\documentclass[11pt]{article}
\usepackage[T1]{fontenc}
\usepackage[utf8]{inputenc}
\usepackage{lmodern,amsmath,amssymb,amsthm,mathtools,microtype}
\usepackage[a4paper,margin=28mm]{geometry}
\PassOptionsToPackage{hyphens}{url}
\usepackage[hidelinks]{hyperref}
\setlength{\emergencystretch}{2.5em}
\newtheorem{theorem}{Theorem}[section]
\newtheorem{proposition}[theorem]{Proposition}
\newtheorem{lemma}[theorem]{Lemma}
\newtheorem{corollary}[theorem]{Corollary}
\theoremstyle{definition}
\newtheorem{definition}[theorem]{Definition}
\newcommand{\workbenchurl}{https://github.com/KokunoYumeto/mathematics-workbenches/blob/main/works/WB-0006.md}
\numberwithin{equation}{section}
\title{The question of Anderson and Ihara on the pro-$p$ Galois action on the projective line minus three
points, at every prime satisfying Vandiver's conjecture}
\author{Entered at deposit}
\date{23 November 2026 \quad CC0 1.0}
\begin{document}
\maketitle
\begin{abstract}
\noindent
Ihara showed that the kernel of the outer action of the absolute Galois group of $\mathbb Q$ on the pro-$p$
fundamental group of the projective line minus three points cuts out a field contained in the maximal pro-$p$
extension of $\mathbb Q(\mu_p)$ unramified outside $p$, and Anderson and Ihara asked whether the two fields are
equal. Sharifi proved that they are equal at odd regular primes if the Deligne--Ihara conjecture holds, and Brown
proved that conjecture. We prove the equality at every odd prime $p$ for which the plus part of the Iwasawa module
of $p$-class groups of the cyclotomic tower is finite. This holds whenever $p$ does not divide the class number of
the real cyclotomic field $\mathbb Q(\mu_p)^+$, which is Vandiver's conjecture at $p$; by the computation of Hart,
Harvey and Ong the equality therefore holds at every odd prime below $2^{31}$. At an irregular prime the Galois
group in question has generators that the regular case does not have. The theorems of Fukaya--Kato and
Wake--Wang-Erickson on Sharifi's symbol map show that cup products of cyclotomic units generate the relevant
second cohomology at every finite level, and we prove that this forces the images of the new generators to have
poles that separate all words. The mathematics and the proofs were produced by GPT-6 Astra (ChatGPT Pro mode).
Claude Opus 5.5 at effort high checked the deduction from Vandiver's conjecture and the present text of Sections 4
to 6 step by step, with the theorems of Fukaya--Kato and Wake--Wang-Erickson read at the places cited; Claude Fable
5.2 at effort max edited the text, wrote its first two sections and refereed Sections 3 to 6. No human has
checked it.
\end{abstract}

\section{The question and the theorem}

Let $p$ be an odd prime,
\begin{equation}
 \mathcal X=\mathbb P^1_{\mathbb Q}\setminus\{0,1,\infty\},\qquad
 \Pi_p=\pi_1^{\mathrm{\acute et}}(\mathcal X_{\overline{\mathbb Q}})^{(p)},
 \qquad
 \rho_p:G_{\mathbb Q}\longrightarrow\operatorname{Out}(\Pi_p),
\end{equation}
where $\Pi_p$ is the maximal pro-$p$ quotient of the geometric fundamental group, a free pro-$p$ group of rank
two, and $\rho_p$ is the outer action of the absolute Galois group. Write $\Omega_p^*$ for the fixed field of
$\ker\rho_p$, and $\Omega_p$ for the maximal pro-$p$ extension of $F_1=\mathbb Q(\mu_p)$ unramified outside $p$,
both inside one algebraic closure. Ihara proved $\Omega_p^*\subseteq\Omega_p$ \cite{Ihara1986}; see
\cite[p.~1]{Sharifi2002}. Anderson and Ihara asked whether
\[
 \Omega_p^*=\Omega_p
\]
\cite[p.~272, question (a)]{AndersonIhara1988}. Put
\begin{equation}
 F_r=\mathbb Q(\mu_{p^r}),\qquad F_\infty=\bigcup_{r\ge1}F_r,
 \qquad X_p=\varprojlim_r\operatorname{Cl}(F_r)\{p\},
\end{equation}
the limit of the $p$-primary parts of the class groups for the norm maps, and let $X_p^+$ be the part of $X_p$
fixed by complex conjugation. Greenberg's conjecture for this tower is the finiteness of $X_p^+$. Vandiver's
conjecture at $p$ is $p\nmid h(F_1^+)$, where $F_1^+$ is the real subfield of $F_1$.

\begin{theorem}\label{thm:greenberg-kernel}
Let $p$ be an odd prime.
\begin{itemize}
\item[\textup{(a)}] If $X_p^+$ is finite, then $\Omega_p^*=\Omega_p$.
\item[\textup{(b)}] If $p$ does not divide the class number of $\mathbb Q(\mu_p)^+$, then $X_p^+=0$, and so
$\Omega_p^*=\Omega_p$.
\item[\textup{(c)}] $\Omega_p^*=\Omega_p$ for every odd prime $p<2^{31}$.
\end{itemize}
\end{theorem}

\paragraph{What it means.}
The pro-$p$ covers of the line minus three points, with the Galois action on them, see every extension of
$\mathbb Q(\mu_p)$ that the ramification condition allows: an element of $G_{\mathbb Q}$ that acts trivially on
the tower of these covers, up to inner automorphisms, fixes $\Omega_p$. Equivalently, the group
$\operatorname{Gal}(\Omega_p/\mathbb Q)$ acts faithfully on $\Pi_p$ by outer automorphisms. Anderson and Ihara
described $\Omega_p^*$ as the field generated by their higher circular $p$-units \cite[Theorem~2\,(IV)]{AndersonIhara1988},
as recalled in \cite[p.~2]{IshiiKernel2026}; for the
primes of the theorem these units therefore generate all of $\Omega_p$. Part (c) uses the verification of
Vandiver's conjecture for every odd prime below $2^{31}$ by Hart, Harvey and Ong
\cite[Section 1, item (2)]{HartHarveyOng2017}; that computation is an input.

\paragraph{Why it matters.}
The action of $G_{\mathbb Q}$ on the fundamental group of $\mathcal X$ is the first case of the programme, going
back to Grothendieck and to Ihara's paper \cite{Ihara1986}, of describing Galois groups through their action on
fundamental groups of moduli spaces. Its pro-$p$ part is the best understood: its image is filtered, the graded
pieces carry the Soul\'e characters, and after tensoring with $\mathbb Q_p$ the graded Lie algebra is free on one
generator in each odd degree at least three, by Brown's theorem. What the kernel is had been known only at regular
primes. The theorem gives it at every prime for which Vandiver's conjecture is known, which at present includes
every irregular prime below $2^{31}$, and it reduces the question at the remaining primes to Greenberg's
conjecture for the cyclotomic tower.

\paragraph{Credit.}
The inclusion $\Omega_p^*\subseteq\Omega_p$ is Ihara's and the question is Anderson and Ihara's. Sharifi proved
that at an odd regular prime the Deligne--Ihara conjecture implies the equality \cite[Theorem 1.1, p.~2]{Sharifi2002},
and Brown proved that conjecture \cite[Theorem 1.1]{Brown2012}; so the regular case is their theorem. The proof
here follows the same plan: embed the arithmetic group in a free prounipotent group and separate its elements by
coefficients. It uses, as theorems of their authors, the weighted completion of Hain and Matsumoto
\cite{HainMatsumoto2001} and Brown's theorem for the target; the theorems of Ferrero--Washington
\cite{FerreroWashington1979} and Iwasawa \cite{Iwasawa1981} for the freeness of the arithmetic group; and, for the
step that is new at irregular primes, the symbol construction of Fukaya and Kato \cite{FukayaKato2012} with the
integral isomorphism theorem of Wake and Wang-Erickson \cite{WakeWangEricksonPseudo2018}, both proved in the
study of Sharifi's conjectures. That Vandiver's conjecture implies Sharifi's conjecture is due to Wake and
Wang-Erickson \cite[abstract and Corollary~5.1.2]{WakeWangErickson2017}. What is supplied here are two lemmas, the comparison of finite-level cup products
with the class-two extension of the arithmetic group (Lemma~\ref{lem:cup-pole-bridge}) and the separation of
Magnus coefficients by poles (Lemma~\ref{lem:rational-magnus}), and the verification of their hypotheses.

\paragraph{Context.}
Sharifi's argument at regular primes uses that $\operatorname{Gal}(\Omega_p/F_1)$ is then free pro-$p$ on
generators in odd characters only. The same paper studies, at irregular primes, whether Soul\'e elements generate
the graded Lie algebra over $\mathbb Z_p$, which is a different question \cite[Sections 1 and 4]{Sharifi2002};
the present theorem makes no integral statement of that kind. Ishii recalls the regular case and proves the
analogue of Sharifi's theorem for once-punctured elliptic curves with complex multiplication, under an analogue of
the Deligne--Ihara conjecture \cite[Theorem 1.1 and Remark 1.1.1, p.~2]{IshiiKernel2026}. In 2019 Sharifi
described the irregular case as ``a very interesting open question of Anderson and Ihara''
\cite[\S1, p.~222]{Sharifi2019}. We have found no earlier statement of the equality at an irregular prime; searches of the literature and of Hexagon on 23 November 2026
returned none.

\paragraph{Scope.}
The theorem does not prove Vandiver's or Greenberg's conjecture and says nothing at an odd prime where $X_p^+$ is
not known to be finite. The prime $2$, where the real place needs a convention, is not treated in this note.

\section{Proof overview}

Let $H=\operatorname{Gal}(\Omega_p/F_\infty)$. An element of $\ker\rho_p$ acts trivially on $\mu_{p^\infty}$, so
its image in $\operatorname{Gal}(\Omega_p/\mathbb Q)$ lies in $H$, and the theorem says that the outer action of
$H$ on $\Pi_p$ is faithful. The action on $\Pi_p$ induces one on its rational unipotent completion, and by Hain
and Matsumoto the image of $H$ there is Zariski dense in a prounipotent group $\mathcal U_p$ whose Lie algebra is,
by Brown's theorem, free on letters $s_m$ of odd weight $m\ge3$ (Section~\ref{sec:rational-image}). So it suffices
to prove that $\rho:H\to\mathcal U_p(\mathbb Q_p)$ is injective.

\emph{The regular case.} $H$ is a free pro-$p$ group (Iwasawa, with $\mu=0$), with an action of
$\Gamma=\operatorname{Gal}(F_\infty/F_1)$ and of $\Delta=\operatorname{Gal}(F_1/\mathbb Q)$. If $p$ is regular,
$H^{\mathrm{ab}}$ is free of rank one over the Iwasawa algebra in each odd character of $\Delta$ and zero in the
even ones. The coefficient of a word in the letters $s_m$ is then an iterated integral of bounded measures on
$\Gamma$, evaluated at the weights $u^m$; a nontrivial element has a nonzero coefficient because a nonzero bounded
measure has a nonzero moment at infinitely many weights.

\emph{The difficulty.} At an irregular prime $H^{\mathrm{ab}}=U\oplus T$ has an even part $T$, of finite rank over
$\mathbb Z_p$, dual to the minus part of $X_p$. The target has no letters of even weight. A generator $y$ of $H$
lying over $T$ is sent to words of length at least two, and nothing in the rational picture prevents the images
of such generators from satisfying relations.

\emph{The idea.} Let $q_e$ be the characteristic polynomial of a generator $\gamma$ of $\Gamma$ on a character
component $T_e$. The element $q_e(\gamma)y$ is trivial in $H^{\mathrm{ab}}$, hence a product of commutators, and its
class modulo triple commutators is an alternating tensor $C_e$ in two odd variables. The length-two coefficient
of $y$ at weights $(m,n)$ is then $C_e(u^m-1,u^n-1)/q_e(u^{m+n})$: it is a quotient of power series with a pole
along $(1+T_1)(1+T_2)=\alpha$ for each root $\alpha$ of $q_e$. If the numerator does not vanish identically on
each of these pole divisors, then in the coefficient of a first nonzero degree one can isolate, by expanding along
the poles of highest order, a single product of nonzero factors; this proves injectivity
(Lemma~\ref{lem:rational-magnus}). Repeated roots are allowed.

\emph{Why the numerators do not vanish.} If $C_e$ vanished on the pole divisor of $\alpha$, the class-two quotient
of $H$ would split equivariantly over an eigenline $W\subset T$ for $\alpha$. A cochain computation then shows
that, at every finite level $n$, the cup product of two even classes in $H^1$ of
$\operatorname{Gal}(\Omega_p/F_{n+1})$ with coefficients $\mathbb Z/p^n(1)$ has zero component along $W$ in
$H^2$ with coefficients $\mathbb Z/p^n(2)$ (Lemma~\ref{lem:cup-pole-bridge}). This is impossible if such cups
generate the even part of that $H^2$ at every level. That generation is the arithmetic input: under the
finiteness of $X_p^+$, Wake and Wang-Erickson prove that Sharifi's symbol map is surjective, and Fukaya and Kato
deduce that the symbols $\{1-\zeta^a,1-\zeta^b\}$ of cyclotomic units generate at each finite level
(Section~6).

\emph{Where the hypothesis is used.} The finiteness of $X_p^+$ is used three times in Section~6: for the
freeness of rank one of $U$ after inverting $p$, for the cyclicity of each $T_e$ after inverting $p$, and for the
surjectivity of the symbol map. Section~7 proves part (b): if $p\nmid h(F_1^+)$ then the $p$-class groups of all
layers of the real tower vanish, because each layer is cyclic of degree $p$ over the previous one and totally
ramified at one prime (Lemma~\ref{lem:persistence}).

\section{The rational outer image}
\label{sec:rational-image}

We first identify the representation into which the arithmetic group
will be embedded. Rational unipotent completion of $\Pi_p$ is
functorial for automorphisms and inner automorphisms. It therefore
induces a homomorphism from its pro-$p$ outer automorphism group to
the outer automorphism group of the rational completion
\cite[Appendix A.6, A.9--A.10]{HainMatsumoto2001}.
In particular, an element acting trivially on $\Pi_p$ acts trivially
on this rational completion.

Let $B_p$ denote the proalgebraic closure of the induced arithmetic
outer image, and let $\mathcal U_p$ denote its prounipotent radical.
The cited Hain--Matsumoto results give
\begin{equation}
 1\longrightarrow\mathcal U_p\longrightarrow B_p
   \longrightarrow\mathbb G_m\longrightarrow1,
 \qquad H\longrightarrow\mathcal U_p(\mathbb Q_p)
\end{equation}
with dense image in each finite weight quotient. The reductive
quotient is the cyclotomic character. Brown's theorem \cite{Brown2012} identifies
the graded Lie algebra of $\mathcal U_p$ with the free Lie algebra
on one letter $s_m$ in each odd arithmetic weight $m\ge3$.

The dense subgroup in Hain--Matsumoto's Theorem~8.4 is the
cyclotomic kernel $I_p^1G_{\mathbb Q,\{p\}}$. To identify its image
here, put $\mathcal G_p=\operatorname{Gal}(\Omega_p/\mathbb Q)$.
The quotient $G_{\mathbb Q,\{p\}}\to\mathcal G_p$ maps that kernel
onto $H$: any lift of an element of $H$ already fixes every
$p$-power root of unity, since the cyclotomic field lies in
$\Omega_p$. Conversely a cyclotomically trivial element maps into
$H=\operatorname{Gal}(\Omega_p/F_\infty)$.
The rational outer action factors through this quotient, so its
images of these two kernels coincide. This proves the asserted
density for the specific arithmetic group used below.

Here is the completion and normalization used below. A Levi splitting
provides the weight decomposition. Choose homogeneous lifts of the
free graded generators. The homomorphism from the completed free Lie
algebra on these lifts induces an isomorphism on every graded piece.
Induction on the finite weight filtration makes it an isomorphism
on each finite quotient; the separated inverse limit is an isomorphism
of completed Lie algebras. Its completed enveloping algebra is thus
the free associative algebra completed by arithmetic weight.
There are only finitely many letters and words below each fixed weight.

For odd $p$, choose a lift $\gamma$ of a generator of
$\operatorname{Gal}(F_\infty/F_1)$ and write
$u=\chi_p(\gamma)\in1+p\mathbb Z_p$. In each positive weight the
linear operator $u^m-1$ is invertible over $\mathbb Q_p$. Successive
conjugation in those weight spaces puts the image of $\gamma$ in
a Levi torus. If the prime-to-$p$ group $\Delta$ commutes with
$\gamma$, its unipotent part is then trivial: a first nonzero weight
of an element centralizing $\gamma$ would be fixed by $u^m$, which
is impossible. Consequently conjugation on the letters is
\begin{equation}
 \gamma s_m=u^m s_m,\qquad
 \delta s_m=\omega(\delta)^m s_m.
\end{equation}
This normalization is performed at finite
weights before taking the inverse limit.

\section{Cup products and the class-two extension}

This section proves the precise link from finite arithmetic cup products
to the denominators that occur in the rational action. All cohomology
of profinite groups is continuous cohomology. Completed tensor and
exterior products are taken in the category of compact
$\mathbb Z_p$-modules.
\begin{lemma}[Finite-level cup comparison]\label{lem:cup-pole-bridge}
Let $p$ be odd, let $P=H\rtimes\Gamma$, where $H$ is free pro-$p$
and $\Gamma=\langle\gamma\rangle\simeq\mathbb Z_p$, and let an
involution $c$ preserve $H$ and fix $\gamma$. Put
\begin{equation}
 \Lambda=\mathbb Z_p[[\Gamma]],\qquad R=\Lambda[1/p],\qquad
 H^{\rm ab}=U\oplus T,
\end{equation}
where $c=-1$ on $U$ and $c=1$ on $T$. Suppose that $U$ is finitely
generated over $\Lambda$, that $U[1/p]$ is free of finite rank over
$R$, and that $T$ is finite free over $\mathbb Z_p$. Suppose also
that each even character component $T_e[1/p]$ under the commuting
prime-to-$p$ action is cyclic over $R$. Write $q_e(z)$ for its monic
characteristic and minimal polynomial for $\gamma$.

Let the coefficient character satisfy $\chi(\gamma)=u\in1+p\mathbb Z_p$
and $\chi(c)=-1$, with $H$ acting trivially on its twists. Assume that,
for every $n\geq1$, the even part of
\begin{equation}
 H^2(P_n,\mathbb Z/p^n(2)),\qquad
 P_n=H\rtimes\Gamma^{p^n},
\end{equation}
is generated by cups of $c$-even classes in
$H^1(P_n,\mathbb Z/p^n(1))$. A cofinal diagonal of the two level
indices suffices.

Then the odd--odd companion numerators of the group extension satisfy the
pole-nonvanishing condition: for every root $\alpha$ of every
nonzero $q_e$, at least one coefficient has nonzero restriction to
\begin{equation}
 D_\alpha=(1+T_1)(1+T_2)-\alpha=0.
                                                               \label{eq:ce1}
\end{equation}
Here coefficients are taken in the bounded-denominator power-series
coordinates of an $R$-basis of $U[1/p]$. Repeated roots are allowed.
For the subsequent pole-separation application the arithmetic module
input supplies rank one in every odd character.
\end{lemma}

\begin{proof}
Since $2$ is a unit, the class-two quotient of $H$ has its integral
Baker--Campbell--Hausdorff Lie structure. Quotient its central
commutator module by all brackets having a $T$ input. The remaining
central module is
\begin{equation}
 N=\widehat{\bigwedge}^{\,2}_{\mathbb Z_p}U.
\end{equation}
Complex conjugation is $1$ on $N$. Thus the minus part of the
resulting Lie module maps isomorphically to $U$, providing a
canonical $\Gamma$-equivariant odd section. Its plus part is the
continuous module extension
\begin{equation}
 0\longrightarrow N\longrightarrow E\longrightarrow T
 \longrightarrow0.                                          \label{eq:ce2}
\end{equation}
Choose a rational cyclic generator $v$ of $T_e[1/p]$ and a lift
$y\in E[1/p]$ in the same character component, using the
prime-to-$p$ character projector. The element
\begin{equation}
 q_e(\gamma)y=C_e\in N[1/p]                                 \label{eq:ce3}
\end{equation}
is precisely the odd--odd quadratic companion error. In the quotient
just taken, brackets involving even lifts vanish, so an ordered
companion product becomes its additive characteristic-polynomial
expression. A derived element has the same degree-two logarithmic
and Magnus coefficient.

Finite generation clears the denominators of the chosen isomorphism
$U[1/p]\simeq R^r$ and its inverse: for some $a$,
$p^a\Lambda^r\subset U\subset p^{-a}\Lambda^r$ in these
coordinates. Completed tensor and exterior squares inherit fixed
denominator bounds. Consequently $N[1/p]$ is the bounded-denominator
alternating two-variable coefficient module, rather than a module of
series with arbitrary coefficient denominators.

Pass to a finite splitting field $L/\mathbb Q_p$, with integers
$\mathcal O$, and suppose that every coefficient of $C_e$ vanishes
on $D_\alpha=0$. All roots satisfy $\alpha-1\in\mathfrak m_{\mathcal O}$,
because the action of the pro-$p$ group $\Gamma$ on the finite-rank
lattice $T$ is continuous. The change of variable
\begin{equation}
 S=(1+T_1)(1+T_2)-\alpha,\qquad
 T_2=\frac{\alpha+S}{1+T_1}-1
\end{equation}
is an integral formal-coordinate change. Its specialization kernel
is $(S)$, and division preserves a common denominator bound.
Since $D_\alpha$ is symmetric, an alternating numerator has an
alternating quotient. Thus, in $N_L$,
\begin{equation}
 C_e=(\gamma-\alpha)A.                                      \label{eq:ce4}
\end{equation}
Wrong-character coefficients are identically zero, so it suffices
to assume vanishing for the matching odd-character pairs.

Set $r_e(z)=q_e(z)/(z-\alpha)$. Cyclicity implies that
$w=r_e(\gamma)v$ is nonzero and spans an $\alpha$-eigenline
$W\subset T_L$. For a repeated root it is the socle of the cyclic
generalized eigenspace. Equations \eqref{eq:ce3}--\eqref{eq:ce4} show that
\begin{equation}
 \widetilde w=r_e(\gamma)y-A
\end{equation}
is a $\Gamma$-equivariant lift of $w$. Hence \eqref{eq:ce2} splits over
$W$ after inverting $p$.

We now compare this split with the cup map of this group. The lattice
$T_{\mathcal O}\cap W$ is saturated, hence an $\mathcal O$-direct
summand of $T_{\mathcal O}$. Choose an $\mathcal O$-linear section
$s:T_{\mathcal O}\to E_{\mathcal O}[1/p]$ extending the equivariant
section on that lattice. A fixed power $p^b$ makes this section
integral. Write the class-two logarithm of $h\in H$ as
\begin{equation}
 x(h)+s(t(h))+z(h),\qquad
 x(h)\in U_{\mathcal O},\quad t(h)\in T_{\mathcal O},\quad
 z(h)\in N_{\mathcal O}[1/p].
\end{equation}
The same $p^b$ clears every $z(h)$. The section defect is
\begin{equation}
 \beta_\delta(t)=\delta s(t)-s(\delta t),\qquad\delta\in\Gamma,
                                                               \label{eq:ce5}
\end{equation}
and it vanishes on $W$.

Fix $n$ and two even cohomology classes $f,g$. Average their cocycles
by $(1+c)/2$. The resulting invariant cocycles vanish on the chosen
$\Gamma_n=\Gamma^{p^n}$ section: $c$ fixes its elements and is
$-1$ on the coefficient module. On $H$ they are additive maps
from $U$ to $\mathcal O/p^n(1)$, vanish on $T$, and are
$\Gamma_n$-equivariant. Put
\begin{equation}
 \psi=f\wedge g:N_{\mathcal O}\longrightarrow\mathcal O/p^n(2),
 \quad
 \psi([x_1,x_2])=f(x_1)g(x_2)-f(x_2)g(x_1).
\end{equation}
The continuous cochain
\begin{equation}
 k(h)=-\frac{p^b}{2}f(h)g(h)-\psi(p^b z(h))                    \label{eq:ce6}
\end{equation}
is integral. The class-two product has central coordinate
$z_1+z_2+[x_1,x_2]/2$. Its substitution in \eqref{eq:ce6} gives
\begin{equation}
 dk(h_1,h_2)=p^b f(h_1)g(h_2).
\end{equation}
Indeed the two summands contribute respectively
$p^b(f_1g_2+f_2g_1)/2$ and $p^b(f_1g_2-f_2g_1)/2$.

Extend $k$ to $P_n$ by $k(h\delta)=k(h)$. For $h,h'\in H$
and $\delta,\delta'\in\Gamma_n$, the residual cocycle is explicitly
\begin{equation}
 \begin{aligned}
 &(p^b f\smile g-dk)(h\delta,h'\delta')\\
 &\qquad=k(\delta h'\delta^{-1})-\chi(\delta)^2k(h')\\
 &\qquad=-\psi\bigl(p^b\beta_\delta(t(h'))\bigr).
 \end{aligned}
\end{equation}
Indeed $f$ and $g$ vanish on the section, their restrictions to $H$
are $\Gamma_n$-equivariant, and the quadratic summand of $k$
transforms by $\chi^2$. Thus its Hochschild--Serre edge is, up to
the convention for the edge sign and input conjugation,
\begin{equation}
 \pm\psi(p^b\beta_\delta(t)).                               \label{eq:ce7}
\end{equation}
This vanishes on $W$. A convention that conjugates $t$ makes no
difference because $W$ is $\Gamma$-stable. Thus the same fixed
$p^b$ multiple of every cup has zero restriction to the dual of
$W$, at every finite level.

For completeness, the target and its transitions are explicit.
Both $H$ and $\Gamma_n$ have $p$-cohomological dimension one,
and the minus part $U$ contributes only odd parity. Therefore
\begin{equation}
 H^2(P_n,\mathcal O/p^n(2))^+
 \simeq H^1(\Gamma_n,\operatorname{Hom}(T,\mathcal O/p^n)(2)).
                                                               \label{eq:ce8}
\end{equation}
Let $V=\operatorname{Hom}_{\mathbb Z_p}(T,\mathcal O)(2)$, with
$\gamma$ operator $B$. Normalized cocycle evaluation identifies
the right side with
\begin{equation}
 V\big/\bigl(p^nV+(B^{p^n}-1)V\bigr).                         \label{eq:ce9}
\end{equation}
Corestriction for the procyclic index-$p$ subgroup sends its
generator value to the same value in the larger group's
coinvariants. Thus the transitions are quotient maps, without an
extra factor of $p$. Continuity implies $B^{p^n}-1\to0$; the
denominators in \eqref{eq:ce9} are cofinal with the $p$-adic neighbourhoods.
Their inverse limit is $V$.

Restriction to the primitive lattice in $W$ gives, after inverting
$p$, a nonzero quotient
\begin{equation}
 V_L\longrightarrow W^*(2),\qquad
 \gamma|_{W^*(2)}=u^2/\alpha,
\end{equation}
of character $\chi^2e^{-1}$, or exponent $2-e$ in cyclotomic
character notation. The restriction maps on the integral Hom
modules are surjective, and the corresponding finite coinvariant
quotients form a separated inverse system. By \eqref{eq:ce7} and the
finite-level cup-surjectivity hypothesis, the image of $p^bV$ in
every one of these quotients is zero. Its image in $W^*(2)$ must
therefore be zero. This contradicts the nonzero torsionfree
quotient. No compatible choice of finite-level cup expressions is
needed, since $b$ was fixed before their selection. The contradiction
proves \eqref{eq:ce1} for every root, including repeated ones.
\end{proof}

\section{Separation in the completed Magnus algebra}

The following lemma upgrades the nonzero quadratic restrictions to
injectivity on the entire free pro-$p$ group. Its coefficient ring is
$\mathbb Z_p[[T]][1/p]$: each series has one common finite denominator.
This bound permits evaluation at the cyclotomic weights and is preserved
by all coordinate changes in a fixed assigned degree.
\begin{lemma}[Separation of Magnus coefficients]\label{lem:rational-magnus}
Let $p$ be odd and let $H$ be a second-countable free pro-$p$ group
with commuting continuous actions of
$\Gamma=\langle\gamma\rangle\simeq\mathbb Z_p$ and
$\Delta\simeq\mathbb F_p^\times$. Let $c\in\Delta$ have order two,
and put
\begin{equation}
 \Lambda=\mathbb Z_p[[\Gamma]],\qquad R=\Lambda[1/p],\qquad
 M=H^{\rm ab}=U\oplus T,
\end{equation}
where $c=-1$ on $U$ and $c=1$ on $T$. Assume that $M$ is finitely
generated over $\Lambda$, that $U[1/p]$ is free of rank one over $R$
in every odd $\Delta$-character, and that $T$ is finite free over
$\mathbb Z_p$, with every nonzero $T_e[1/p]$ cyclic over $R$.
Assume also that abelian eigenclasses have exact $\Delta$-eigenlifts,
including inside arbitrarily small $\Delta$-stable neighbourhoods.

Let $\rho:H\to\mathcal U(\mathbb Q_p)$ be continuous and
$\Gamma\times\Delta$-equivariant, where $\mathcal U$ is the
weight-completed free prounipotent group on letters $s_m$ in all odd
weights $m\ge3$. Suppose its image is Zariski dense in every finite
weight quotient and, after a Levi normalization,
\begin{equation}
 \gamma s_m=u^m s_m,\qquad \delta s_m=\omega(\delta)^m s_m,
\end{equation}
where $u$ topologically generates $1+p\mathbb Z_p$.

For each nonzero $T_e$, choose an integral multiple $v_e$ of a rational
cyclic generator, an exact eigenlift $y_e\in H$, and its
conjugates $y_{e,j}=\gamma^j y_e\gamma^{-j}$. Let $q_e(z)$ be the
characteristic polynomial of $\gamma$ on $T_e$. In a rational
$R$-basis of $U[1/p]$, let $C_{e;a,b}(T_1,T_2)$ be the odd--odd
quadratic coefficients of the companion error formed from these conjugates. Suppose that
for every root $\alpha$ of every $q_e$, there is a matching ordered
odd pair $(a,b)$ for which
\begin{equation}
 C_{e;a,b}\left(T,\frac{\alpha}{1+T}-1\right)\ne0.
                                                               \label{eq:gm1}
\end{equation}
Then $\rho$ is injective. All coefficient series here have a common
finite denominator in each fixed assigned degree; no bound uniform
over all degrees is required.
\end{lemma}

\begin{proof}
\emph{A compact free basis.}
A compact torsionfree $\mathbb Z_p$-module is a product of copies of
$\mathbb Z_p$: its Pontryagin dual is a divisible $p$-primary group,
hence a direct sum of copies of $\mathbb Q_p/\mathbb Z_p$.
The abelianization of free $H$ is torsionfree and metrizable.
Using the integral character idempotents, choose pointed convergent
$\mathbb Z_p$-bases in the $U_a$, and finite bases in the $T_e$.
The open quotient map $H\to M$ permits lifts of the convergent
basis vectors tending to $1$. Apply the exact-eigenlift construction
inside a nested sequence of $\Delta$-stable open neighbourhoods;
it preserves both the abelian classes and their convergence.

Let $F$ be the free pro-$p$ group on this pointed sequence and the
finite even set. Its abelianization is the product module on those
basis vectors, so the resulting map $F\to H$ is an isomorphism on
abelianizations. Its compact image is Frattini-surjective and hence
equals $H$. Indeed a proper closed subgroup of any pro-$p$ group is
contained in a proper open subgroup and then in an index-$p$ subgroup,
so cannot be Frattini-surjective. Since $H$ is free, the surjection
has a continuous section. The section is inverse on abelianizations;
its closed image is Frattini-surjective in $F$ and is all of $F$.
Thus $F\simeq H$. This argument does not use an infinite-rank
Hopficity assertion.

Give the odd Magnus variables $X_i$ assigned degree $1$, and the
finite even variables $Y_l$ assigned degree $2$. This Magnus
expansion is injective. To see this directly, kill a tail of the
convergent free sequence. Every finite quotient kills such a tail,
so $H$ is the inverse limit of these finite-rank free groups.
An ordinary finite-rank Magnus expansion detects every nonidentity
element in some finite assigned degree. At each fixed degree the
coefficient module is a finite sum of completed tensor products of
$U$ and $T$, using these assigned degrees.

\emph{Odd rational coordinates.}
Choose a character-preserving isomorphism
\begin{equation}
 U[1/p]\simeq\bigoplus_{a\ {\rm odd}}Re_a.
\end{equation}
Finite generation over $\Lambda$ clears the denominators of this
isomorphism and its inverse. Thus there are continuous integral
maps between $U$ and $\Lambda^{(p-1)/2}$ whose two compositions are
multiplication by the same power of $p$. Their completed tensor
powers have compositions equal to the corresponding powers of that
scalar. The tensor modules are torsionfree: in the product bases
they are products of copies of $\mathbb Z_p$. Hence the maps are
injective and become inverse isomorphisms after inverting $p$.
Every fixed degree therefore has coordinates in
\begin{equation}
 \mathbb Z_p[[\Gamma^s]][1/p]
   =\mathbb Z_p[[T_1,\ldots,T_s]][1/p],                       \label{eq:gm2}
\end{equation}
for its singleton positions and their character types. These are
the Amice transforms of bounded measures on $\Gamma^s$.
The denominator bound may depend on the degree, which causes no
problem at a first nonzero degree.

The target linear coefficient of weight $m$ is a continuous additive
map from $M$ to $\mathbb Q_p$. Parity kills $T$, and character
equivariance kills every $U_a$ except $a\equiv m\pmod{p-1}$.
On the remaining rank-one $R$-module, $\Gamma$-equivariance makes
it a scalar multiple of evaluation at $T=u^m-1$: check this on
polynomials in $\gamma$ and then pass to bounded series by
continuity. The scalar is nonzero by Zariski density, since the
weight-$m$ abelian coordinate of the target cannot vanish on a
dense image. Rescale $s_m$ by this scalar. Such weightwise rescaling
is an automorphism of the completed free Lie algebra, because each
finite weight quotient has finitely many generators. The singleton
kernel is now exactly evaluation at $u^m-1$ in its matching type.

\emph{Even companion coordinates.}
Put $\lambda_e=\operatorname{rank}_{\mathbb Z_p}T_e$. The classes
\begin{equation}
 v_e,\gamma v_e,\ldots,\gamma^{\lambda_e-1}v_e
\end{equation}
are a rational basis, though they need not be an integral basis.
The polynomial $q_e$ is monic integral, has unit constant term,
and reduces to $(z-1)^{\lambda_e}$ modulo $p$, since the image of
$\Gamma$ in $\operatorname{GL}(T_e/p)$ is a finite $p$-group.
It is also the minimal polynomial, by rational cyclicity.

The group conjugates $y_{e,j}$ have assigned degree at least $2$
and target generator length at least $2$. Their assigned-degree-two
parts have the form
\begin{equation}
 Z_{e,j}=\sum_l b_{j,l}Y_{e,l}+A_{e,j}(X),                    \label{eq:gm3}
\end{equation}
where $b$ is a finite invertible rational matrix and $A_{e,j}$ is
a completed quadratic odd tensor. Before the odd coordinate change
these coefficients are integral; afterwards they retain bounded
denominators. There are no other terms of assigned degree two.
The quadratic tensor is alternating, by the commutator description
or primitiveness of the first Magnus coefficient.

Replace the finite variables $Y$ by $Z$. This is a homogeneous
triangular automorphism of the completed coefficient algebra:
it fixes $X$ and its inverse is $Y=b^{-1}(Z-A(X))$.
In each fixed degree only finitely many even variables occur,
so both directions preserve a common finite denominator bound.
In particular a first nonzero degree remains nonzero.
This is a change in the rational tensor algebra, not a claim that
the companion classes form an integral group basis.

By definition of $Z_{e,j}$, its leading target is the length-two
part of the group element $y_{e,j}$. For an ordered target pair
$(m,n)$, put $z=u^{m+n}$. Its coefficient satisfies
\begin{equation}
 B_{e,j;m,n}=z^jB_{e,0;m,n}.
\end{equation}
Write $q_e(z)=z^{\lambda_e}-\sum_jc_{e,j}z^j$ and form
\begin{equation}
 E_e=(\gamma^{\lambda_e}y_e\gamma^{-\lambda_e})
       \left(\prod_jy_{e,j}^{c_{e,j}}\right)^{-1}.
\end{equation}
Its abelian class is zero, so $E_e$ is in the closed derived
subgroup. Its assigned-degree-two part is the odd alternating
tensor $C_e$. There are no target length-one terms in any even
factor, so comparison of length two gives exactly
\begin{equation}
 q_e(z)B_{e,0;m,n}
   =C_{e;a,b}(u^m-1,u^n-1),\qquad
 a\equiv m,\ b\equiv n\pmod{p-1}.                           \label{eq:gm4}
\end{equation}
This is the companion class of the extension in
Lemma~\ref{lem:cup-pole-bridge}. Exact eigenlifts make its assigned
degree-two character $e$; hence wrong-character pairs have zero
coefficient. One can also obtain this zero from \eqref{eq:gm4} on the full
residue grid and the successive analytic identity theorem.

\emph{Partial fractions and continuity.}
Pass to a finite splitting field $L$, with integers $\mathcal O$.
All roots $\alpha$ satisfy $\alpha-1\in\mathfrak m_{\mathcal O}$.
The finite-dimensional rational-function space with basis
$z^j/q_e(z)$ has alternative basis
\begin{equation}
 (z-\alpha)^{-k},\qquad
 1\le k\le\operatorname{mult}_e(\alpha).
\end{equation}
The corresponding constant invertible change in the finite $Z$
variables preserves assigned degree and bounded denominators in
every fixed degree. Its even letter $(e,\alpha,k)$ has kernel
\begin{equation}
 \frac{C_{e;a,b}(T_1,T_2)}
      {((1+T_1)(1+T_2)-\alpha)^k}.                          \label{eq:gm5}
\end{equation}
The numerator is common to all letters from the same companion
component, even when roots repeat.

There are finitely many positive integers $w$ such that $u^w$ is a
root of some $q_e$. Restrict all test weights to a sufficiently
large ordinary tail, excluding every vanishing pair denominator.
In each fixed odd residue class this still includes sequences
$m_0+(p-1)p^j$ with an interior $p$-adic accumulation point, for
a sufficiently large fixed representative $m_0$.

Substitution is legitimate at every finite arithmetic-weight cutoff.
In the compact free basis, truncate its convergent sequence and
apply finite-rank Magnus substitution. These truncations converge
to the identity on $H$, since every finite group quotient kills a
tail. The image of $H$ in the finite target quotient is compact,
so its coefficient functions are bounded and continuous. Passing
to the limit therefore gives the completed tensor substitution.
The bounded rational coordinate changes preserve this statement in
each fixed degree. In \eqref{eq:gm2} it is equivalently integration against
bounded measures and the finite Fubini rule. We do not invoke a
free pro-$p$ universal property for a noncompact rational group.

\emph{Separation of a first nonzero degree.}
Suppose $g\ne1$ and $\rho(g)=1$, and let $d$ be its first nonzero
assigned degree after these invertible changes. A larger source
degree cannot contribute to target length $d$. Its coefficient is
a finite sum over ordered tilings of $d$ positions by singleton
odd types and pairs labelled $(e,\alpha,k)$:
\begin{equation}
 \begin{gathered}
 \sum_{\mathcal P}A_{\mathcal P}(T_i:\ i\text{ singleton})
       \prod_{(j,e,\alpha,k)\in\mathcal P}
       \frac{C_{e;a_j,a_{j+1}}(T_j,T_{j+1})}
            {D_{j,\alpha}^{k}}=0,\\
 D_{j,\alpha}=(1+T_j)(1+T_{j+1})-\alpha.
 \end{gathered}                                            \label{eq:gm6}
\end{equation}
Every $A_{\mathcal P}$ has a common bounded denominator. Clear
the finitely many denominators. Vanishing on the Cartesian grid
of large odd test weights and the successive one-variable identity
theorem make \eqref{eq:gm6} an identity in the fraction field of
$\mathcal O[[T_1,\ldots,T_d]][1/p]$.

Select a nonzero pattern with the largest number of pairs. Fix its
pair positions $J$ and its roots and even characters, and among
nonzero patterns with those data choose a componentwise maximal
pole-order vector $(k_j)$. At each selected pair choose the ordered
odd types in \eqref{eq:gm1}, and at each singleton its source type. These
choices are independent because the pairs are disjoint.

Set $S_j=D_{j,\alpha_j}$ for $j\in J$. Retain as base variables
the singleton variables and the left variable in each selected pair.
The substitutions
\begin{equation}
 T_{j+1}=\frac{\alpha_j+S_j}{1+T_j}-1
\end{equation}
are automorphisms of the complete integral power-series ring.
Let $K_0$ be the fraction field of the bounded-denominator base
ring. The identity embeds in $K_0[[S_j:j\in J]][S_j^{-1}:j\in J]$,
with finite lower bounds on exponents. All other denominators are
units in $K_0[[S_j]]$: a different root at the same position has
nonzero constant root difference; a different pair position joins
two distinct selected blocks, whose surviving variables are
independent, so its product is not identically a prescribed root.
Wrong even characters have already vanished.

Take the coefficient of $\prod_{j\in J}S_j^{-k_j}$.
A pattern missing a selected position or using a different root
there has no negative exponent in that variable. A pattern
containing all selected pairs and an extra pair has zero coefficient
by the first maximality choice. Among remaining patterns, a
contribution requires all pole orders to dominate $k_j$; the
componentwise maximal choice excludes strict domination and an
incomparable vector fails in at least one variable. Singleton types
must match. Thus exactly one coefficient remains:
\begin{equation}
 A_{\mathcal P}(T_i:\ i\text{ singleton})
 \prod_{j\in J}
 C_{e_j;a_j,b_j}\left(T_j,\frac{\alpha_j}{1+T_j}-1\right).
\end{equation}
The first factor is nonzero and its coefficient ring injects into
$K_0$; every other factor is nonzero by \eqref{eq:gm1}. Their product is
nonzero in a field, a contradiction. If $J$ is empty, the same
argument is the ordinary singleton moment identity. This proves
injectivity of $\rho$.
\end{proof}


\section{The arithmetic consequence of Greenberg's hypothesis}

We now verify all hypotheses of the two preceding lemmas for the
cyclotomic Galois group. The following comparison puts the finite-level
symbol theorem in the group-cohomological form required there.
\begin{lemma}[Low-degree \'etale comparison]\label{lem:etale-prop}
Let $S$ be a connected normal affine scheme, let
$G=\pi_1^{\mathrm{\acute et}}(S)$, and let $P$ be its maximal
pro-$p$ quotient. For a finite $p$-primary local coefficient module
$A$ whose action factors through $P$, the natural maps give
\begin{equation}
 H^1(P,A)\simeq H^1_{\mathrm{\acute et}}(S,A),\qquad
 H^2(P,A)\hookrightarrow H^2_{\mathrm{\acute et}}(S,A).
\end{equation}
They are compatible with cup products and with finite external
automorphisms. Consequently generation of a character summand of
the target $H^2$ by cups of $H^1$ classes gives the same generation
in group cohomology and makes the displayed injection an
isomorphism on that summand.
\end{lemma}

\begin{proof}
Put $N=\ker(G\to P)$. An element of $H^1(N,A)^P$ is a continuous
$P$-equivariant homomorphism $N\to A$, since $N$ acts trivially
on $A$. If its image were nonzero, its kernel would be $G$-normal
and the quotient of $G$ by that kernel would be an extension of
$P$ by a nonzero finite $p$-group. This is a pro-$p$ quotient:
every finite quotient has $p$-power order. It would be strictly
larger than the maximal pro-$p$ quotient, a contradiction.
Thus $H^1(N,A)^P=0$. The five-term Hochschild--Serre sequence gives
\begin{equation}
 H^1(P,A)\simeq H^1(G,A),\qquad
 H^2(P,A)\hookrightarrow H^2(G,A).
\end{equation}
No finite-generation hypothesis on $N$ was used.

The full universal finite \'etale tower over $S$ has zero $H^1$
with finite locally constant coefficients: after trivializing the
coefficients, an $H^1$ torsor is finite \'etale and becomes trivial
at a later stage of the same tower. Finite \'etale objects and
torsors descend under these affine limits. The Cartan--Leray
spectral sequence for the tower therefore gives
\begin{equation}
 H^1(G,A)\simeq H^1_{\mathrm{\acute et}}(S,A),\qquad
 H^2(G,A)\hookrightarrow H^2_{\mathrm{\acute et}}(S,A).
\end{equation}
This is only its low-degree consequence, and asserts no vanishing
of the higher rows. The sequence, the profinite-tower passage and
the simply connected $H^1$ vanishing are given explicitly in
\cite[Theorem~14.9, Remark~14.10, Example~14.11 and
Proposition~14.13, pp.~96--98]{MilneLEC2013}.
All comparison maps are natural and respect cups.
Lift each cup generator using the $H^1$ isomorphism and add its
group cups; this proves surjectivity onto the asserted summand.
Injectivity then identifies every group class in that summand with
the lifted sum.
\end{proof}

\begin{proof}[Proof of the odd-prime assertion of Theorem~\ref{thm:greenberg-kernel}]
At $p=3$, Sharifi's Theorem 1.1 and Brown's theorem apply, since $3$ is regular. Hence
assume $p\ge5$. Put
\begin{equation}
 P=\operatorname{Gal}(\Omega_p/F_1),\quad
 H=\operatorname{Gal}(\Omega_p/F_\infty),\quad
 \Gamma=\operatorname{Gal}(F_\infty/F_1),\quad
 M=H^{\rm ab}=U\oplus T.
\end{equation}
Choose the finite $\Delta$ lift by the coprime complement theorem.
The fibre over a generator of $\Gamma$ is a $\Delta$-stable
$H$-torsor because the cyclotomic quotient is abelian.
The coprime identity $H^1(\Delta,H)=1$, obtained from the finite
complement theorem and compactness, gives a $\Delta$-fixed lift
$\gamma$. Its procyclic closure maps isomorphically to $\Gamma$
and supplies the commuting section used above.
Ferrero--Washington gives $\mu=0$ for this abelian cyclotomic tower
\cite{FerreroWashington1979}. Iwasawa's Theorem 2 then states that
$H$ is free pro-$p$ \cite[p.~273]{Iwasawa1981}.
The maximal extension in that theorem is the extension $\Omega_p$
used here: a finite $p$-extension over $F_\infty$ descends to a finite
cyclotomic layer, and the normal closure of its descended
$p$-extension over $F_1$ is again a $p$-extension unramified outside
$p$.
The finite-level Kummer sequence makes $H^1(P,\mathbb F_p)$
finite dimensional, so $P$ is finitely generated and $H$ is
second-countable. Adjust a finite generating set of $P$ by powers
of the chosen $\Gamma$ section to obtain elements of $H$.
Together with $\gamma$ they generate $P$; their closed normal
closure is $H$, since the remaining quotient is the given
procyclic quotient. Their abelian classes therefore generate $M$
over $\Lambda=\mathbb Z_p[[\Gamma]]$.
Section~\ref{sec:rational-image} gives the continuous dense rational
representation and its Levi normalization as required in
Lemma~\ref{lem:rational-magnus}. The exact eigenlift construction
is \cite[Lemma~4, p.~3 of arXiv:math/0104116v2]{Sharifi2002}; performing it inside
stable neighbourhoods preserves convergence. We verify the module and cup hypotheses.

\emph{Rational odd module.}
The ring $R=\Lambda[1/p]$ is a PID. Indeed $\Lambda$ is a
two-dimensional regular local UFD, and its only height-two prime
contains $p$. Thus $R$ is a one-dimensional noetherian UFD and
hence a PID. It remains the ring of bounded-denominator series,
not $\mathbb Q_p[[T]]$.

For an odd character $a\ne\omega,\omega^{-1}$, apply the full
duality calculation in Wake--Wang-Erickson's Proposition~2.2.6
with $\theta=\omega a^{-1}$. It gives
\begin{equation}
 \operatorname{Ext}^1_\Lambda(M_a^{\#}(1),\Lambda)=X_\theta,
 \qquad
 \operatorname{Ext}^2_\Lambda(M_a^{\#}(1),\Lambda)=0.           \label{eq:gm7}
\end{equation}
Here $M$ is their $p$-ramified module $\mathfrak X$, whereas $X$
on the right is unramified. Since $\theta$ is even, Greenberg
makes the right-hand side zero after inverting $p$. Ext commutes
with this localization for finitely generated modules. Over a PID,
the Ext group into the ring of a finitely generated torsion summand
is nonzero unless that summand is zero. Thus $M_a[1/p]$ is free.
Its rank is one by the local/global unit exact sequence in the
same proof: local units have rank one in this character, global
minus units contribute only roots of unity, and the unramified
quotient is torsion. The inversion $\#$ and Tate twist transport
the ring by automorphisms, so undoing them preserves freeness
\cite[Proposition~2.2.6 and proof, pp.~3--4]{WakeWangErickson2017}.

The two excluded characters can be calculated directly.
Coleman's exact sequence at tame level one is
\begin{equation}
 0\longrightarrow\mathbb Z_p(1)\longrightarrow U_{\rm loc}
 \longrightarrow\mathbb Z_p[[\Delta\times\Gamma]]
 \longrightarrow\mathbb Z_p(1)\longrightarrow0
\end{equation}
\cite[Theorem~5.5.31, p.~148; Proposition~6.1.1,
pp.~151--152]{SharifiIwasawaNotes}.
Global minus units are exactly the root-of-unity term, and
Stickelberger gives $X_\omega=X_{\omega^{-1}}=0$; the latter
exceptional statement is explicitly recorded in
\cite[proof of Lemma~5.2.7, p.~13]{WakeWangErickson2017}.
The $\omega^{-1}$ component of $M$ is consequently free of rank
one. In the $\omega$ component, quotienting the local units by
their root-of-unity kernel identifies $M_\omega$ with
\begin{equation}
 \ker(\Lambda\longrightarrow\mathbb Z_p(1))
     =(\gamma-u)\Lambda,
\end{equation}
again free of rank one. The two characters are distinct for
$p\ge5$. The valuation direction at the unique prime above $p$
is even and does not affect this calculation. We have therefore
proved the required rational freeness for every odd character.

\emph{Rational even module.}
Kummer duality identifies $T=M^+$ up to finite kernel and cokernel
with $(X^-)^\iota(1)$
\cite[Proposition~3.4.7 and Corollary~3.4.8,
p.~92]{SharifiIwasawaNotes}.
Ferrero--Washington makes it finitely generated over $\mathbb Z_p$,
and freeness of $H$ excludes $\mathbb Z_p$-torsion. Thus $T$ is
finite free. Greenberg implies that $X^-$ is pseudocyclic
character by character, by the exact local/global unit sequence and
Iwasawa adjunction in
\cite[Corollary~2.1.3 and Conjecture~2.1.5 with the following
explanation, p.~8]{WakeWangEricksonPseudo2018}.
Finite kernels and cokernels vanish after inverting $p$, so every
$T_e[1/p]$ is cyclic. This allows repeated roots; it does not assert
integral cyclicity without Vandiver.

\emph{Integral symbol surjectivity under Greenberg.}
Fix an even character $\theta\ne1,\omega^2$ and put
$\chi=\omega^{-1}\theta$. Wake--Wang-Erickson's
Corollary~8.1.3(2) assumes exactly
\begin{equation}
 \begin{gathered}
 X_\chi(1)[1/p]\text{ cyclic over }(\Lambda/\xi_\chi)[1/p],\\
 X_\theta\otimes_\Lambda(\Lambda/\xi_\chi)[1/p]=0.
 \end{gathered}                                            \label{eq:gm8}
\end{equation}
The localization in the second condition is essential.
Both conditions follow from the just stated consequences of
Greenberg. The conclusion is an integral isomorphism
\begin{equation}
 \varpi:(\mathcal H^-_\theta/I_\theta\mathcal H^-_\theta)
           /(p\text{-primary torsion})
       \xrightarrow{\ \sim\ }X_\chi(1),                    \label{eq:gm9}
\end{equation}
so the original $\varpi$ is surjective
\cite[Corollary~8.1.3 and complete proof,
pp.~48--49]{WakeWangEricksonPseudo2018}.
The source proves this with repeated roots: the derivative identity
makes the error nilpotent in each cyclic local factor, rather than
requiring the derivative to be invertible. Thus \eqref{eq:gm9} does not
insert a simple-root hypothesis.

Fukaya--Kato, \S7.1.10, deduce from this surjectivity that for every
$r\ge1$ the finite-level $\theta$ component is generated by the
symbols
\begin{equation}
 \{1-\zeta_{p^r}^{a},1-\zeta_{p^r}^{b}\}.
\end{equation}
Their \S7.1.9 identifies its reduction modulo $p^r$ with the
corresponding class-group quotient, and their \S6.1.5 shows that
at tame level one the only excluded even characters are
$1,\omega^2$
\cite[\S6.1.5, p.~76; \S\S7.1.8--7.1.10,
p.~86]{FukayaKato2012}.
Those excluded characters correspond to the unramified components
$X_{\omega^{-1}}$ and $X_\omega$, which are zero. Norm projections
from $X^-$ onto each finite minus class group are surjective, so
the omitted finite targets are zero as well
\cite[Proposition~3.4.1 and proof,
pp.~90--91]{SharifiIwasawaNotes}.
Components with unit characteristic power series are also zero:
the source's $X^-$ is $p$-torsionfree, and a finite such module must
vanish. We apply the nonunit running assumption of the modular
source only to the nonzero components. Hence every even target is
covered.

\emph{All finite group levels and coefficient conventions.}
Let
\begin{equation}
 P_n=H\rtimes\Gamma^{p^n},\qquad
 S_n=\operatorname{Spec}\mathcal O_{F_{n+1}}[1/p].
\end{equation}
The maximal pro-$p$ quotient of $\pi_1(S_n)$ is exactly $P_n$.
Indeed the normal closure over $F_1$ of any finite Galois
$p$-extension of $F_{n+1}$ unramified outside $p$ remains a
$p$-extension: its conjugate compositum over $F_{n+1}$ has
$p$-group Galois group, and $F_{n+1}/F_1$ has $p$-power degree.
Thus it lies in the fixed field $\Omega_p$ used throughout.

At $F_{n+1}$, the finite-level symbol theorem gives generation
with coefficients $p^{n+1}$. Reduction gives generation with
coefficients $p^n$, exactly as in Lemma~\ref{lem:cup-pole-bridge}.
The class-group quotient description in FK \S7.1.9 makes this
reduction explicit. The Brauer contribution is zero: these fields
have no real places, and only one finite prime is omitted, so the
global invariant condition forces its lone possible local
invariant to vanish.

For a $p$-unit $v$, the quotient $v/c(v)$ is a global unit of
absolute value one at every embedding and hence is a root of
unity. Consequently the negative $c$ part of the $p$-unit group
modulo $p^n$ is cyclic. Its alternating self-cup is zero because
$2$ is a unit; mixed-parity cups lie in the odd target. Thus the
even target is generated by cups of $c$-even unit classes. This
argument does not identify the whole of $H^1$ with units and does
not use Vandiver.

Apply Lemma~\ref{lem:etale-prop}. Each Kummer class has a unique
lift to $H^1(P_n,\mathbb Z/p^n(1))$, and group $H^2$ injects into
\'etale $H^2$. Lifting the even symbol generators makes this
injection surjective on the even part and gives precisely
\begin{equation}
 H^2(P_n,\mathbb Z/p^n(2))^+
  =\bigl\langle f\smile g\mid
       f,g\in H^1(P_n,\mathbb Z/p^n(1))^+\bigr\rangle.       \label{eq:gm10}
\end{equation}
This treats every $n\ge1$; no compatibility of expressions as sums
of cups is needed. Scalar extension to the integers of a finite
splitting field preserves the assertion, since that ring is finite
free over $\mathbb Z_p$.

Lemma~\ref{lem:cup-pole-bridge} now gives the nonzero companion
restriction \eqref{eq:gm1} for every root of every even component, including
repeated roots. Its dual root is $u^2/\alpha$ and its dual character
is $\omega^{2-e}$, so the inversion and both Tate twists match
the symbol target. Lemma~\ref{lem:rational-magnus} proves
injectivity on the entire $H$.

Finally an element in the kernel of the original pro-$p$ tripod
action is cyclotomically trivial, hence lies in $H$. Functorial
rational completion sends it to $1$ under $\rho$, and the proved
injectivity makes it trivial. This is exactly
$\Omega_p^*=\Omega_p$.
\end{proof}
\section{Vandiver and the verified finite range}

We prove directly that Vandiver supplies the Iwasawa hypothesis in
Theorem~\ref{thm:greenberg-kernel}.

\begin{lemma}\label{lem:persistence}
Let $L/K$ be cyclic of odd prime degree $p$, ramified at precisely one
finite prime, which is totally ramified. If $\operatorname{Cl}(K)\{p\}=0$,
then $\operatorname{Cl}(L)\{p\}=0$.
\end{lemma}
\begin{proof}
Let $E/L$ be the Hilbert $p$-class field. Its uniqueness makes $E/K$
Galois, and $G=\operatorname{Gal}(E/K)$ is a finite $p$-group. The
inertia group $I$ above the ramified prime has order $p$, since
$E/L$ is unramified. Let $N$ be its normal closure in $G$. All other
inertia groups at that prime are conjugate, and every other prime
is unramified; an odd-degree extension introduces no infinite-place
ramification. Thus $E^N/K$ is everywhere unramified. If $G/N$ were
nontrivial, it would have a cyclic quotient of order $p$, contradicting
class field theory and the hypothesis on $K$. Hence $N=G$.
The Frattini quotient of $G$ is consequently generated by the image
of one generator of $I$, so $G$ is cyclic. Its subgroup $I$ is then
normal, whence $G=N=I$ and $E=L$.
\end{proof}

At every step $F_{r+1}^+/F_r^+$ the sole ramified finite prime is the
prime above $p$, and it is totally ramified. Lemma~\ref{lem:persistence}
therefore gives $\operatorname{Cl}(F_r^+)\{p\}=0$ at all layers when
Vandiver holds at $F_1$. Extension and norm of ideals identify
$\operatorname{Cl}(F_r^+)\{p\}$ with
$\operatorname{Cl}(F_r)\{p\}^+$: their composites are multiplication
by $2$ and $1+c$, and $2$ is invertible on a $p$-group. Taking norm
limits gives $X_p^+=0$. The odd-prime assertion already proved now
establishes the Vandiver part of Theorem~\ref{thm:greenberg-kernel}.

Hart--Harvey--Ong verify Vandiver for every odd prime below $2^{31}$
\cite[Section 1, item (2)]{HartHarveyOng2017}. Their computation
proves the stated numerical
corollary. No assertion about an unverified prime is needed for that
finite range.


\section{What remains open}

The theorem determines the kernel field at every prime for which its
arithmetic hypothesis has been established. For an odd prime with
$X_p^+$ finite, the proof passes through the following specific maps:
the modular symbol map to the cyclotomic class module; its projections
to finite-level cup products; the edge map for
$1\to H\to H\rtimes\Gamma^{p^n}\to\Gamma^{p^n}\to1$;
and the coefficients of the induced rational outer representation.
The two Tate twists in the edge map send an eigenvalue $\alpha$ to
$u^2/\alpha$. This identifies the arithmetic direction tested by each
pole, including generalized eigenspaces for repeated roots.

To establish the same conclusion for every prime by this argument one
must establish the required arithmetic module structure and cup
nonvanishing beyond the finiteness hypothesis. The proof here does not
provide that step. Vanishing of a particular quadratic coefficient
would also require further analysis: a group element can have a
nonzero image at higher generator length. Thus no converse to the
criterion is asserted.

There is a related integral question concerning generation of the
weight-graded Lie algebra by specified Soul\'e elements. The map from
a group to its weight-graded Lie algebra, followed by tensoring with
$\mathbb Q_p$, can discard integral indices. Our separation argument
uses nonzero rational leading coefficients; it never proves that those
coefficients are $p$-adic units. Accordingly the kernel-field theorem
makes no integral generation assertion. Sharifi treats the maps from
free integral generators to the filtered group and graded Lie algebra
explicitly in \cite[Sections 1 and 4]{Sharifi2002}.

\section*{Provenance and checks}

The mathematics, and the text of Sections 3 to 8, are the work of GPT-6 Astra (ChatGPT Pro mode), in research
runs of November 2026; the same model reread its proofs and its applications of the cited theorems. Sections 1
and 2 and this note were written by Claude (Claude Fable 5.2 at effort max), as editor, who also removed the
treatment of the prime $2$ from the text and changed it in no other mathematical respect.

Claude Opus 5.5 at effort high checked every step of Lemma~\ref{lem:persistence} and of the deduction of
$X_p^+=0$ from it in the present text (23 November). In earlier versions of the argument it checked the
separation lemma and the deduction of the theorem at primes satisfying Vandiver's conjecture (21 and 22
November). On 23 November it then checked the present text of Sections 4 to 6 step by step, with
\cite[Proposition~2.2.6]{WakeWangErickson2017}, \cite[Corollary~2.1.3, Conjecture~2.1.5 and
Corollary~8.1.3(2)]{WakeWangEricksonPseudo2018} and \cite[\S6.1.5 and \S\S7.1.9--7.1.10]{FukayaKato2012} opened at
those places, and found no mathematical error; the statements cited from \cite{MilneLEC2013},
\cite{SharifiIwasawaNotes} and \cite{Iwasawa1981} were not opened again. The editor, Claude Fable 5.2 at effort max, then refereed
Sections 3 to 6 against the cited theorems of Hain--Matsumoto, Wake--Wang-Erickson and Fukaya--Kato and found no error. Claude Opus 5.5 at effort max read the whole note and checked Section 7 in full. No human has checked the proofs or this text.

The sources were opened by the run, which recorded for each the pages used; the locators in the text are its
own. For this note the editor read the abstracts of \cite{Sharifi2002} and \cite{IshiiKernel2026}; the literature
check of 23 November opened \cite{Sharifi2002}, \cite{Sharifi2019}, \cite{IshiiKernel2026}, Sharifi's 2011
paper on his conjectures, \cite{FukayaKato2012} and both papers of Wake and Wang-Erickson at the source. What
Section 1 says about \cite{AndersonIhara1988} follows \cite[p.~2]{IshiiKernel2026}, Ihara's inclusion follows
\cite[p.~1]{Sharifi2002}, and the sentences on the programme are the editor's summary.

The records of the checks are at the permanent page of this work, WB-0006:

{\raggedright\smallskip\noindent\url{\workbenchurl}.\par}

\smallskip
This text is dedicated to the public domain (CC0 1.0).

\begingroup
\small
\begin{thebibliography}{99}
\bibitem{AndersonIhara1988}
G.~Anderson and Y.~Ihara, \emph{Pro-$l$ branched coverings of $\mathbb P^1$ and higher circular $l$-units},
Ann. of Math. (2) \textbf{128} (1988), 271--293.
\url{https://doi.org/10.2307/1971443}.
The kernel question is on p.~272, question (a); its statement is also reproduced in
\cite[p.~2]{IshiiKernel2026}.

\bibitem{Brown2012}F.~Brown, \emph{Mixed Tate motives over $\mathbb Z$},
Ann. of Math. \textbf{175} (2012), 949--976.
\url{https://annals.math.princeton.edu/wp-content/uploads/annals-v175-n2-p10-p.pdf}.

\bibitem{FerreroWashington1979}
B.~Ferrero and L.~C.~Washington,
\emph{The Iwasawa invariant $\mu_p$ vanishes for abelian number fields},
Ann. of Math. (2) \textbf{109} (1979), 377--395.
\url{https://doi.org/10.2307/1971116}.

\bibitem{FukayaKato2012}T.~Fukaya and K.~Kato,
\emph{On conjectures of Sharifi}, author preprint,
15 August 2012, 121 pages.
\url{https://www.math.ucla.edu/~sharifi/sharificonj.pdf}.
Locators in this author version:
\S\S6.1.4--6.1.8, pp.~76--77;
\S7.1.3, pp.~84--85;
Conjecture~7.1.8 and \S\S7.1.9--7.1.10, p.~86;
Theorem~7.2.8 and \S\S7.2.10--7.2.12, p.~89;
Lemma~10.2.2 and \S10.2.3, p.~107.

\bibitem{HainMatsumoto2001}R.~Hain and M.~Matsumoto, \emph{Weighted completion of Galois groups
and Galois actions on the fundamental group of $\mathbb P^1-\{0,1,\infty\}$},
arXiv:math/0006158v2, 11 August 2001.
\url{https://arxiv.org/pdf/math/0006158}.

\bibitem{HartHarveyOng2017}W. B. Hart, D. Harvey and W. Ong,
Irregular primes to two billion,
\emph{Math. Comp.} 86 (2017), 3031--3049.
\url{https://doi.org/10.1090/mcom/3211}.
Author preprint:
\url{https://arxiv.org/abs/1605.02398}.
Section 1, item (2), verifies Vandiver for every odd prime
\(\ell<2^{31}\).

\bibitem{Ihara1986}
Y.~Ihara, \emph{Profinite braid groups, Galois representations and complex multiplications},
Ann. of Math. (2) \textbf{123} (1986), 43--106.
\url{https://doi.org/10.2307/1971352}.

\bibitem{IshiiKernel2026}S.~Ishii, \emph{On the kernels of the pro-$p$ outer Galois representations
associated to once-punctured CM elliptic curves},
arXiv:2312.04196v3, 6 January 2026.
\url{https://arxiv.org/pdf/2312.04196}.

\bibitem{Iwasawa1981}K.~Iwasawa, \emph{Riemann--Hurwitz formula and $p$-adic Galois
representations for number fields},
Tohoku Math. J. \textbf{33} (1981), 263--288.
\url{https://www.jstage.jst.go.jp/article/tmj1949/33/2/33_2_263/_pdf/-char/en}.

\bibitem{MilneLEC2013}J.~S.~Milne,
\emph{Lectures on \textup{\'Etale} Cohomology},
version~2.21, 22 March 2013, 202 pages.
\url{https://www.jmilne.org/math/CourseNotes/LEC.pdf}.
Theorem~14.9, Remark~14.10, Example~14.11 and Proposition~14.13,
printed pp.~96--98.

\bibitem{Sharifi2002}R.~T.~Sharifi, \emph{Relationships between conjectures on the structure
of pro-$p$ Galois groups unramified outside $p$},
Proc. Sympos. Pure Math. \textbf{70} (2002), 275--284.
Author's 13-page version: \url{https://www.math.ucla.edu/~sharifi/publ/galstr.pdf};
arXiv:math/0104116v2.

\bibitem{Sharifi2019}R.~T.~Sharifi, \emph{Soul\'e characters in the work of Ihara}, RIMS K\^oky\^uroku
\textbf{2120} (2019), 221--228.
\url{https://www.kurims.kyoto-u.ac.jp/~kyodo/kokyuroku/contents/pdf/2120-11.pdf}.

\bibitem{SharifiIwasawaNotes}
R.~T.~Sharifi, \emph{Iwasawa Theory}, author lecture notes, 184 pp.,
retrieved 23 November 2026.
\url{https://www.math.ucla.edu/~sharifi/notes/iwasawa.pdf}.
Locators: Propositions 3.4.1 and 3.4.7, Corollary 3.4.8, pp.~90--92;
Theorem 5.5.31, p.~148; Proposition 6.1.1, pp.~151--152.

\bibitem{WakeWangErickson2017}P.~Wake and C.~Wang-Erickson,
\emph{Ordinary pseudorepresentations and modular forms},
Proc. Amer. Math. Soc. Ser.~B \textbf{4} (2017), 53--71.
\url{https://doi.org/10.1090/bproc/29}.
Final author preprint arXiv:1510.01661v4, 30 January 2017:
\url{https://arxiv.org/pdf/1510.01661v4}.
Locators in that version: \S2.1, pp.~2--3;
Proposition~2.2.6, p.~4; Theorem~4.2.8, pp.~10--11;
Corollaries~5.1.1--5.1.2, p.~11; Lemma~5.2.7, p.~13.

\bibitem{WakeWangEricksonPseudo2018}P.~Wake and C.~Wang-Erickson,
\emph{Pseudo-modularity and Iwasawa theory},
Amer. J. Math. \textbf{140} (2018), no.~4, 977--1040.
\url{https://doi.org/10.1353/ajm.2018.0022}.
Author version: arXiv:1505.05128v3,
7 September 2017, 54 pages:
\url{https://arxiv.org/pdf/1505.05128v3}.
Locators in that version: Corollary~2.1.3 and Conjecture~2.1.5, p.~8;
Corollaries~8.1.2--8.1.3, pp.~48--49;
Theorem~8.2.3, pp.~50--51.
\end{thebibliography}
\endgroup
\end{document}
